Neste capítulo são apresentados os conceitos fundamentais e o levantamento bibliográfico para a compreensão do problema da pesquisa. São discutidos os conceitos sobre os paradigmas relacional e não relacional (NoSQL), comparação técnica entre os formatos \texttt{JSON} e \texttt{JSONB} no PostgreSQL e as especificidades dos bancos de dados PostgreSQL e MongoDB. Por fim, apresenta uma análise dos principais trabalhos que servem como base para a solução desenvolvida, contextualizando como estudos anteriores abordaram comparações entre esses sistemas e identificando as lacunas que o presente trabalho se propõe a preencher.

\section{Conceitos e Técnicas Relevantes}
\subsection{Bancos Relacionais e Bancos Não Relacionais}
\label{sec:relacional-vs-nao-relacional}

A evolução dos sistemas de armazenamento de dados é marcada pela busca por eficiência, integridade e escalabilidade diante do volume crescente e da variedade de informações. Assim, duas principais abordagens atendem requisitos de aplicações modernas: bancos de dados relacionais e não relacionais, logo, compreender as características de cada paradigma é essencial para a escolha arquitetural em sistemas como catálogos dinâmicos.

Sistemas Gerenciadores de Bancos de Dados Relacionais (SGBDR) organizam os dados em estruturas tabulares com esquema predefinido, 
o que favorece o tratamento de dados estruturados \cite{matcha2025database}. Nesses sistemas, a confiabilidade transacional é sustentada pelas 
propriedades ACID: atomicidade, consistência, isolamento e 
durabilidade, que asseguram a execução íntegra das transações,  
a preservação de invariantes do sistema, a não interferência entre 
transações concorrentes e a persistência dos efeitos após falhas 
\cite{waterpeace2026sok}.
Por essa razão, bancos relacionais 
permanecem amplamente utilizados em domínios críticos e permitem consultas complexas envolvendo múltiplas tabelas.

Bancos de dados não relacionais (NoSQL) surgiram como resposta às limitações dos sistemas relacionais tradicionais diante das crescentes demandas de armazenamento e consulta de grandes volumes de dados \cite{davoudian2018}. O paradigma NoSQL abrange diferentes categorias conforme a organização dos dados, como chave-valor, documentos, bancos orientados a colunas e grafos. Entre essas, os bancos orientados a documentos destacam-se por permitir que atributos dinâmicos e heterogêneos coexistam no mesmo documento sem estruturas previamente fixadas (\textit{schema-on-read}). Contudo, a literatura não trata o NoSQL como substituto universal do modelo relacional, mas como uma alternativa complementar \cite{nayak2013}.

\subsection{JSON vs JSONB}

O PostgreSQL oferece dois tipos de dados nativos para o armazenamento de documentos no formato JSON: 
JSON e JSONB. JSON armazena dados em formato textual, preservando a representação original 
do documento inserido. Em contraste, JSONB armazena dados em formato binário, processando o 
documento durante a inserção para otimizar seu armazenamento e acesso \cite{Liu2020}.

A escolha por JSONB em relação a JSON é fortemente justificada pelas implicações de 
desempenho. Embora JSONB demande processamento adicional durante a inserção para realizar 
a conversão binária, esse custo é compensado significativamente nas operações de leitura e 
indexação. 

Estudos indicam que o tipo de dado JSONB, por armazenar documentos em formato 
binário, oferece melhor desempenho em consultas, com tempos de execução 
menores em diferentes cenários avaliados \cite{Aji2026}, e permite indexação eficiente via 
índices GIN, enquanto o formato textual de JSON requer processamento custoso durante operações 
de consulta \cite{Liu2020}. Para catálogos dinâmicos com padrões de acesso centrados em leitura, 
essa vantagem de desempenho é crucial.


\subsection{PostgreSQL com JSONB}

O PostgreSQL é um sistema gerenciador de banco de dados relacional de código aberto
que utiliza a linguagem SQL (\textit{Structured Query Language}) para definição,
manipulação e consulta dos dados. Os dados são armazenados em
tabelas organizadas em páginas de 8~KB e para permitir acesso
concorrente, adota o controle de concorrência multiversão (MVCC)\footnote{\url{https://www.postgresql.org/docs/current/mvcc-intro.html}}: leitores consultam
versões anteriores dos registros enquanto escritores atualizam a versão corrente, de modo
que um não bloqueia o outro \cite{Bondugula2024}. Para acelerar as consultas, o PostgreSQL oferece diferentes tipos de
índice, como B-tree\footnote{\url{https://www.postgresql.org/docs/current/btree.html}}, \textit{Generalized Search Tree} (GiST)\footnote{\url{https://www.postgresql.org/docs/current/gist.html}}, GIN e \textit{Block Range Index} (BRIN)\footnote{\url{https://www.postgresql.org/docs/current/brin.html}}.

Nesse contexto, o tipo JSONB faz com que o PostgreSQL seja híbrido, onde atributos estáveis
permanecem em colunas relacionais e atributos variáveis, como os de um catálogo
dinâmico, ficam em uma coluna JSONB. A abordagem proposta busca combinar o desenvolvimento rápido com JSON com a consistência do modelo relacional e a eficiência do SQL
\cite{Irshad2019}.

Neste trabalho utiliza-se três tipos de índice:

\begin{itemize}
    \item \textbf{Índice simples:} um índice B-tree sobre uma única coluna. É a estrutura
    padrão do PostgreSQL, com busca de complexidade logarítmica, e apresenta bom
    desempenho em consultas de seleção e atualização.
    \item \textbf{Índice composto:} um índice B-tree sobre múltiplas colunas, indicado
    para filtros que combinam vários atributos. A ordem das colunas
    importa, pois o índice só é plenamente aproveitado por consultas que filtram suas
    colunas iniciais.
    \item \textbf{Índice GIN (\textit{Generalized Inverted Index}):} um índice invertido
    que mapeia cada chave e valor presentes nos documentos aos registros que os contêm,
    sendo o indicado para colunas JSONB.
\end{itemize}

\subsection{MongoDB}

O MongoDB é um banco de dados não relacional orientado a documentos de código aberto, projetado para oferecer alta disponibilidade, escalabilidade horizontal 
e flexibilidade de esquema. Em vez de tabelas e linhas, o MongoDB organiza os dados em \textit{coleções} compostas por \textit{documentos}. Internamente, os documentos 
são armazenados no formato BSON (\textit{Binary JSON})\footnote{\url{https://www.mongodb.com/docs/manual/reference/bson-types/}}, uma representação binária do JSON que organiza o documento como uma lista ordenada de elementos,
cada um composto por nome do campo, tipo de dado e valor \cite{Truica2013}.

Para o gerenciamento de armazenamento no disco, o MongoDB utiliza por padrão o motor
\textit{WiredTiger}\footnote{\url{https://www.mongodb.com/docs/manual/core/wiredtiger/}} \cite{Nguyen2016}. Esse motor emprega um mecanismo de concorrência no nível de documento
com controle otimista \cite{MongoDBConcurrency}, permitindo que múltiplas operações de
leitura e escrita ocorram simultaneamente na mesma coleção sem bloquear a coleção inteira.
Assim como no PostgreSQL, a eficiência no acesso aos documentos e a execução de consultas no MongoDB dependem diretamente do uso adequado de índices, 
destacando-se quatro tipos principais utilizados neste trabalho:

\begin{itemize}
    \item \textbf{Índice de Campo Único (\textit{Single Field Index})}\footnote{\url{https://www.mongodb.com/docs/manual/core/indexes/index-types/index-single/}}\textbf{:} estrutura de busca baseada em árvore sobre um único campo de nível superior ou sobre 
    um caminho específico dentro de um documento, com busca de complexidade logarítmica.
    
    \item \textbf{Índice Composto (\textit{Compound Index})}\footnote{\url{https://www.mongodb.com/docs/manual/core/indexes/index-types/index-compound/}}\textbf{:} estrutura de busca baseada em árvore sobre múltiplos campos de um mesmo documento. 
    A sequência dos campos na definição do índice é fundamental, assim como no índice composto do PostgreSQL.
    
    \item \textbf{Índice Chave-Múltipla (\textit{Multikey Index})}\footnote{\url{https://www.mongodb.com/docs/manual/core/indexes/index-types/index-multikey/}}\textbf{:} estrutura de busca baseada em árvore, acionada automaticamente pelo MongoDB quando o campo indexado contém um \textit{array}, 
    permitindo que consultas filtrem com eficiência documentos que possuam valor dentro de coleções ou chaves aninhadas.

    \item \textbf{Índice Curinga (\textit{Wildcard Index})}\footnote{\url{https://www.mongodb.com/docs/manual/core/indexes/index-types/index-wildcard/}}\textbf{:} indexa dinamicamente todos os campos e subdocumentos presentes na coleção (ou em um caminho específico). É a estrutura genérica indicada para cenários com atributos altamente 
    heterogêneos e imprevisíveis.
\end{itemize}


\section{Trabalhos Relacionados}

A literatura recente investiga de forma ampla o desempenho comparativo entre sistemas de gerenciamento de bancos de dados relacionais e NoSQL. Para fundamentar a análise proposta neste trabalho, destacam-se os seguintes estudos relevantes.

\citeonline{Ortiz2016} conduziram um estudo comparativo avaliando o desempenho do PostgreSQL com a introdução do tipo \texttt{JSONB} frente ao MongoDB em cargas de até 4 milhões de documentos. Analisando métricas de tamanho do banco de dados, tempo de carga e tempo de resposta em consultas de seleção, os autores observaram que o PostgreSQL apresentou tempos de inserção e armazenamento em disco altamente competitivos. Nas consultas de leitura, o MongoDB manteve-se mais rápido em seleções com ordenação, enquanto o PostgreSQL com \texttt{JSONB} superou o MongoDB em operações de agrupamento (\textit{GROUP BY}). O estudo é fundamental por demonstrar a evolução do PostgreSQL em capacidades não relacionais, porém limitou-se a uma comparação genérica sem aprofundar em padrões estruturais específicos de JSONB ou em estratégias de indexação avançadas.

\citeonline{Asiminidis2018b} investigaram as características de desempenho do PostgreSQL ao comparar a estrutura relacional tradicional contra colunas contendo dados semiestruturados nos formatos \texttt{JSON} e \texttt{JSONB}, utilizando um conjunto de dados de IoT com mais de 1 milhão de registros. Os autores observaram que a inserção em campos \texttt{JSON}/\texttt{JSONB} é mais rápida do que no esquema relacional devido à ausência de checagens estritas de tipo na ingestão. Ao comparar os formatos semiestruturados, enquanto o tipo \texttt{JSON} armazena uma cópia exata do texto de entrada, o \texttt{JSONB} armazena os dados em bytes decompostos. Isso permitiu ao \texttt{JSONB} superar o \texttt{JSON} nas seleções e agregações, concluindo que o \texttt{JSONB} representa o melhor meio-termo entre a flexibilidade do \texttt{JSON} e o desempenho de consulta do modelo relacional.

\citeonline{Urnikiene2026} apresentaram uma análise comparativa focada em cenários de \textit{e-commerce}, avaliando o desempenho de leitura do PostgreSQL e do MongoDB sob volumes de até 300 mil produtos. Os resultados demonstraram que o PostgreSQL executou consultas analíticas complexas mais rápido do que o MongoDB, mantendo maior estabilidade e menor desvio padrão nas medições. Por outro lado, o MongoDB apresentou uma eficiência de armazenamento em disco superior devido ao encapsulamento em documentos incorporados. Os autores concluíram que esquemas relacionais normalizados mantêm vantagens expressivas em cenários de \textit{e-commerce} onde cargas analíticas dominam o sistema, ressaltando que a escolha do banco deve ser guiada pela complexidade das consultas e volume de dados.

Estes trabalhos, embora fundamentais para evidenciar a diferença entre SGBDs relacionais e NoSQL, apresentam limitações para o cenário de catálogos dinâmicos 
de grande porte como o do iFood. Entre elas, destacam-se o tratamento genérico do \texttt{JSONB}, sem distinguir como a estrutura interna dos documentos afeta o desempenho,
a exploração limitada da indexação, sem comparar índices simples, compostos e GIN no PostgreSQL com os índices equivalentes do MongoDB e
a ausência de uma análise de como o desempenho varia com a escalabilidade. O presente estudo diferencia-se ao propor uma análise de um recorte do cenário real, 
categorizando os atributos flexíveis do catálogo em três padrões estruturais de \texttt{JSONB}, utilizando diferentes tipos de indexação, além de quantificar o 
comportamento do PostgreSQL e do MongoDB isoladamente sob três eixos de escalabilidade em ambiente de contêineres Docker\footnote{\url{https://docs.docker.com/}}.


\section{Considerações Finais}

Este capítulo estabeleceu a fundamentação teórica necessária, apresentando os paradigmas relacional e não relacional, o tipo \texttt{JSONB} do PostgreSQL e o MongoDB, com seus respectivos tipos de índice. A revisão dos trabalhos relacionados evidenciou a lacuna existente na análise de padrões estruturais de \texttt{JSONB}, de estratégias de indexação e de escalabilidade em catálogos dinâmicos. O Capítulo~\ref{chapter:desenvolvimento} detalhará a metodologia adotada, abrangendo a modelagem e a geração dos dados, as consultas e os índices avaliados, bem como os experimentos de escalabilidade e os resultados obtidos.
